\section{İlgili Çalışmalar}\label{sec:ilgili-calismalar}

\subsection{ILS ve vuruşun yeri}
Double-bridge vuruşu \citet{martin_1992_largestep}'in büyük adımlı Markov zincirleriyle yaygınlaşmış,
\citet{lourenco_2010_ils} ise ILS'yi yerel arama, pertürbasyon ve kabul kuralından oluşan bir çatı
olarak sistematikleştirmiştir. Hızlı 2-opt ve Or-opt uygulamalarının aday listeleri ve bakma
bitleri (don't-look bits) \citep{bentley_1992_fast} ile pertürbasyon sayısının çözüm kalitesindeki
belirleyici rolü \citep{johnson_1997_case} bu çatının pratik temelidir. \citet{applegate_2003_clk}
büyük örneklerde vuruşu turun küçük bir bölgesine sınırlamanın (segment, rastgele yürüyüş,
geometrik vuruşlar) önemini gösterir. LKH \citep{helsgaun_2000_lkh} ve LKH-3 \citep{helsgaun_2017_lkh3}
bu çizginin en güçlü çözücüleridir. Bu çalışmalar vuruşun \emph{nerede} yapılacağını çoğunlukla
rastgele ya da yerellik kuralıyla belirler. Double-bridge seçiminin rastgele yerine yanlı
yapılmasının yararı ILS çatısında daha önce tartışılmıştır \citep{lourenco_2010_ils}; bizim
katkımız, aynı motoru sabit tutup yanlılığın kaynağını öğrenilmiş ve eğitimsiz skorlar arasında eş
bütçede karşılaştırmaktır.

\subsection{Uzun kenar bilgisiyle yönlendirme}
GLS, yerel optimumdaki kenarları $c_e/(1+p_e)$ fayda ölçütüyle cezalandırır: uzun ve az
cezalandırılmış kenarlar önce hedeflenir \citep{voudouris_1999_gls}. Bilgi güdümlü yerel arama (knowledge-guided local search, KGLS) aynı fikri araç rotalamada pertürbasyon aşamasına taşır: bir veri madenciliği çalışmasından türetilen rota genişliği ve kenar uzunluğu ölçütlerine göre ``kötü'' kenarları cezalandırır ve yerel aramayı bu kenarlardan başlatır. Yazarlar bu yönlendirmenin rastgele cezalandırmadan daha iyi sonuç verdiğini, ancak etkisinin yerel arama operatörlerinin seçiminden küçük kaldığını bildirir \citep{arnold_2019_kgls}. Basit
çözüm ve örnek ölçütlerinin iyi ve kötü çözümleri ayırt edebildiği de gösterilmiştir
\citep{arnold_2019_vrpgood}. LKH'nin alfa-yakınlık aday kümeleri ise bir kenarın değerini, o kenar zorunlu kılındığında en küçük 1-ağacın uzunluğundaki artışla ölçer \citep{helsgaun_2000_lkh}. Fazla-uzunluk vekilimiz, GLS'nin uzun kenar sezgisini kenar başına $O(1)$ işlem maliyetiyle
double-bridge pertürbasyon hedeflemesine uyarlar. Katkımız bu yalın uyarlamanın öğrenilmiş
skorlara karşı eş duvar-saatinde ve ön-kayıtlı bir protokolle sınanmasıdır.

\subsection{Öğrenilmiş arama rehberliği}
Uçtan uca inşa yöntemleri turu doğrudan üretir \citep{bello_2017_nco, dai_2017_s2v, kool_2019_attention, kwon_2020_pomo};
iyileştirme yöntemleri hamle seçimini öğrenir \citep{wu_2022_improvement}. Daha güçlü sonuçlar
veren kol, öğrenilmiş bileşeni klasik bir aramanın içine koyar. \citet{hudson_2022_gnngls} GNN'nin
kenar pişmanlık tahminlerini GLS cezalarına bağlar; \citet{sui_2024_neuralgls} bu fikri geliştirir.
NeuroLKH aday kümelerini \citep{xin_2021_neurolkh}, \citet{zheng_2023_rlkh} aynı kararları takviye
öğrenmesiyle öğrenir. Isı haritası yöntemleri kenar olasılıklarını aramaya taşır \citep{fu_2021_generalize, qiu_2022_dimes, sun_2023_difusco}; hiyerarşik ayrıştırma ise öğrenilmiş politikaları çok büyük örneklere ölçekler \citep{pan_2023_htsp}. Taramalar bu yöntemleri öğrenmenin çözücüdeki rolüne (uçtan uca, iyileştirme, klasik sezgiselle melez) ya da metasezgiselin hangi bileşenine (başlatma, operatör seçimi, parametre ayarı) katıldığına göre sınıflar \citep{bengio_2021_mlco, karimimamaghan_2022_mlmeta, sui_2025_survey}; eğitim dağılımı dışına
genellemenin zor olduğu da bilinir \citep{joshi_2020_generalization}. Bizim GNN, üyelik-MLP ve
kazanç-MLP skorlayıcılarımız bu ailenin vuruş hedeflemedeki temsilcileridir. Farkımız, bunları
aynı bilgiden türetilen eğitimsiz vekillerle eş duvar-saatinde karşılaştırmamızdır.

\subsection{Basit tabanlara karşı sınama}
Birkaç çalışma, öğrenilmiş bileşenin katkısını eğitimsiz bir tabana karşı ayrıştırmanın neden
gerekli olduğunu göstermiştir. Maksimum bağımsız küme (maximum independent set, MIS) için GNN güdümlü ağaç aramasında
GNN çıktısı rastgele değerlerle değiştirildiğinde başarım benzer kalmıştır \citep{bother_2022_wrong};
modern GNN'lerin aynı problemde klasik açgözlü algoritmanın gerisinde kaldığı da bildirilmiştir
\citep{angelini_2023_gnngreedy}. Bu sonuçlar MIS'e aittir ve burada yalnız benzetme olarak anılır.
TSP'de \citet{xia_2024_rethinking}, uzaklıktan türetilen basit bir ısı haritasının öğrenilmiş ısı
haritalarıyla rekabet ettiğini gösterir. Rotalama için önerilen hesaplamalı test kılavuzu güçlü
klasik tabanlar ve adil zaman raporu ister \citep{accorsi_2022_guidelines}. Çalışmamız bu deseni
ILS vuruş hedeflemesine, dokunulmamış bir sette ve ön-kayıtlı testlerle taşır.

\subsection{Uyarlanabilir operatör seçiminde kredi ataması}
Reaktif arama, parametreleri arama geçmişinden uyarlar \citep{battiti_2009_reactive}. Uyarlanabilir
operatör seçimi (adaptive operator selection, AOS) çok kollu haydut çerçevesini kullanır
\citep{auer_2002_ucb, dacosta_2008_aos}; uyarlanabilir geniş komşuluk araması da operatörleri geçmiş
başarımlarına göre seçer \citep{ropke_2006_alns}. Bu alanda başarım hem seçim kuralına \citep{thierens_2005_pursuit} hem de operatöre verilen kredinin nasıl tanımlandığına bağlıdır \citep{fialho_2010_aos}. Ortalama ödül nadir ama büyük
iyileşmeleri bastırabilir; uç değerlere dayalı kredi bu nedenle önerilmiştir \citep{fialho_2008_extreme}.
Çevrimiçi vuruş skorlayıcımız bu sınıftandır. Bulgumuz, vuruş başına ortalama ödülün bu motorda
son kaliteyle hizalanmadığını ölçen bir örnektir; çevrimiçi öğrenmenin genel olarak zararlı olduğu
iddiası değildir.

\subsection{Deneysel yöntembilim}
Sezgisel testinin yarışma yerine kontrollü ve hipotez güdümlü bir deney olması gerektiği uzun
süredir savunulur \citep{hooker_1995_testing}. Deneysel algoritma çalışmalarında örnek seçimi,
tekrarlanabilirlik ve raporlama kuralları \citet{johnson_2002_guide}, \citet{mcgeoch_2012_guide} ve
\citet{bartzbeielstein_2020_benchmarking} tarafından derlenmiştir. Ön-kayıt, doğrulayıcı tahmini
sonradan yapılan açıklamadan ayırır \citep{nosek_2018_prereg}. Çok problemli karşılaştırmada örneği
birim alan Wilcoxon testi ve Holm düzeltmesi yerleşik uygulamadır
\citep{demsar_2006_comparisons, wilcoxon_1945, holm_1979}. ``Fark yok'' iddiası ise önceden belirlenmiş
bir sınırla kurulan tek yönlü testler ister \citep{schuirmann_1987_tost, lakens_2017_equivalence}.
Protokolümüz bu öğeleri, geliştirme ve doğrulama setlerini ayıran iki aşamalı bir zincirde birleştirir.
